% ============================================================
%  NLDB 2026 --- Speech notes for the 15-minute talk
%  Paper: Towards Robust Uzbek Neural Dependency Parsing
%         (Cross-Treebank Training & Linguistically-Motivated
%          Subword Fusion)
%  Author: Sanatbek Matlatipov
%
%  Optimised for printing on a black & white printer:
%    - 12pt serif body, 1.4 line spacing
%    - bold cues only (no colours), clear slide separators
%    - generous left margin for handwritten annotations
%    - cumulative timing budget printed next to each slide
%
%  Build with:  pdflatex presentation_speaker_notes.tex
% ============================================================
\documentclass[12pt,a4paper]{article}

\usepackage[utf8]{inputenc}
\usepackage[T1]{fontenc}
\usepackage{lmodern}
\usepackage[margin=2.2cm,left=3.2cm,right=2.2cm]{geometry}
\usepackage{setspace}
\usepackage{enumitem}
\usepackage{titlesec}
\usepackage{fancyhdr}
\usepackage{microtype}
\usepackage{parskip}
\usepackage{amsmath,amssymb}

\setstretch{1.40}
\setlist[itemize]{leftmargin=1.3em,topsep=2pt,itemsep=2pt}

\titleformat{\section}
   {\normalfont\Large\bfseries\filright}
   {Slide \thesection.}{0.6em}{}
\titlespacing*{\section}{0pt}{1.2em}{0.4em}

\pagestyle{fancy}
\fancyhf{}
\fancyhead[L]{\small NLDB 2026 --- Speech notes (15 min)}
\fancyhead[R]{\small S.~Matlatipov}
\fancyfoot[C]{\thepage}
\renewcommand{\headrulewidth}{0.4pt}

\newcommand{\timing}[2]{%
   \par\noindent\textbf{\small [#1\,$\to$\,#2\,min]}\par\vspace{0.2em}}
\newcommand{\say}[1]{\textbf{#1}}
\newcommand{\pause}{\par\vspace{0.4em}\textbf{[pause]}\par}

\begin{document}

% ------------------------------------------------------------
%  COVER PAGE
% ------------------------------------------------------------
\begin{center}
   {\LARGE\bfseries Speech notes --- 15 minutes}\\[0.6em]
   {\large Towards Robust Uzbek Neural\\Dependency Parsing}\\[0.3em]
   {\normalsize Cross-Treebank Training \& Linguistically-Motivated Subword Fusion}\\[0.6em]
   {\normalsize International Conference on Applications of Natural Language\\to Information Systems}\\
   {\normalsize \textbf{NLDB 2026} --- 18 June 2026}\\[1.2em]
   \textit{Sanatbek Matlatipov}\\
   National University of Uzbekistan, Tashkent\\
   \texttt{s.matlatipov@nuu.uz}
\end{center}

\vspace{1em}
\noindent\textbf{How to use these notes.}
Each section maps to one slide. The bracketed times
\textbf{[start\,$\to$\,end]} are cumulative running time. Phrases in
\textbf{bold} are the ones to emphasise; \textbf{[pause]} marks a short
breath. Plain text is the spoken script --- read naturally, not
word-for-word.

\vspace{0.5em}
\noindent\textbf{Total budget:} $\approx$15 minutes of speech inside a
15-minute slot. Target around 0:30--0:45 per slide; the motivation now
spans two slides --- \textbf{Slide 4 (challenges)} and \textbf{Slide 5
(suffixes carry the syntax)} --- so keep each tight. Slow down on
\textbf{Slide 14 (fusion)} and \textbf{Slides 18--19 (results)}.

\vspace{0.5em}
\noindent\textbf{Two ideas the audience must remember:}
\textbf{last-subword fusion} and \textbf{cross-treebank training}.
\textbf{Numbers to land:} 85.08 / 72.39 / 63.81 (UPOS / UAS / LAS) and
cross-treebank \textbf{+9.62 to +11.16 LAS}. Also mention: \textbf{Uzbek
is spoken by 40 million}, \textbf{Turkic Family}, \textbf{SOV word
order}, \textbf{null subject}, \textbf{no articles}.

\vfill\hrule\vspace{0.5em}
\noindent\textit{Print double-sided on A4. The wide left margin is for handwritten cues.}

\newpage

% ------------------------------------------------------------
%  SLIDE 1 --- Title
% ------------------------------------------------------------
\section*{Slide 1 --- Title}
\timing{0:00}{0:20}

Good morning / afternoon, and thank you to the NLDB organisers.

My name is \say{Sanatbek Matlatipov}, from the National University of
Uzbekistan. My talk is on \say{Towards Robust Uzbek Neural Dependency
Parsing} --- making Uzbek parsing robust in a genuinely
\say{low-resource} setting, using two lightweight and efficient ideas
rather than bigger models.

% ------------------------------------------------------------
%  SLIDE 2 --- Outline
% ------------------------------------------------------------
\section*{Slide 2 --- Outline}
\timing{0:20}{0:35}

The talk has a simple arc:

\begin{itemize}
   \item Quick \textbf{motivation} and related work.
   \item The \textbf{dataset} we built and pair with.
   \item The \textbf{methodology} --- pipeline, representation, model.
   \item And \textbf{results}, error analysis, and takeaways.
\end{itemize}

% ------------------------------------------------------------
%  SLIDE 3 --- Uzbek Language I
% ------------------------------------------------------------
\section*{Slide 3 --- Uzbek Language I}
\timing{0:35}{1:05}

Before diving into the parsing challenge, let me briefly introduce Uzbek.

Uzbek NLP has seen \say{growing interest} in recent years ---
TahrirchiBERT, BERTbek, and several shared tasks demonstrate that.
But compared to major world languages, computational resources remain
sparse, and we still face genuinely low-resource conditions for
structured prediction tasks like parsing.

Uzbek is the \say{official language} of Uzbekistan. It is native to
Central Asia, Afghanistan, parts of Russia, and Xinjiang in China.
Roughly \textbf{40 million} people speak Uzbek as of 2026 --- a
significant speech community, yet under-resourced in computational
linguistics.

The map on the right shows where Uzbek is spoken. You can see it spans
a large geographic region across Central Asia and beyond.

% ------------------------------------------------------------
%  SLIDE 4 --- Why is Uzbek dependency parsing hard? (challenges)
% ------------------------------------------------------------
\section*{Slide 4 --- Why is Uzbek dependency parsing hard?}
\timing{1:05}{1:35}

Three things make Uzbek parsing hard. First, Uzbek is \say{agglutinative}
--- a single word stacks suffixes for case, tense, person, even
evidentiality. Look at the figure: \textit{kelmaganlardanmisizlar?} is
\say{one} Uzbek word, but it unpacks into the eight-word English question
``Are you one of those who didn't come?'' --- come, not, past, plural, one
of, question, are, you. That is the morphological density a parser has to
decode. Second, \say{flexible word order} together with strong genre
variation. Third, it is genuinely \say{low-resource} --- fewer than a
thousand annotated UD sentences before this work.

On the right are the \textbf{Turkic-family} traits the parser must cope
with: \say{null subject} --- you can drop the subject pronoun;
\say{no articles} and \say{no grammatical gender}, unlike English; and
\textbf{Subject-Object-Verb} order, quite different from English SVO.

\pause

So the picture is \say{rich morphology} on one side, \say{scarce data} on
the other. Those are exactly the two problems we attack --- better
representations for the morphology, more data for the scarcity. The next
slide shows \say{why} that morphology matters so much.

% ------------------------------------------------------------
%  SLIDE 5 --- Suffixes carry the syntax (the hook)
% ------------------------------------------------------------
\section*{Slide 5 --- Suffixes carry the syntax}
\timing{1:35}{2:05}

Here is the intuition in one example. Take \textit{bolalarning} --- ``of
the children.'' Stem \textit{bola}, plus plural \textit{-lar}, plus
genitive \textbf{\textit{-ning}}. That final suffix \textbf{is} the
syntactic relation: the \textbf{\textit{-ning}} genitive is what creates
the possessive arc --- \texttt{nmod:poss} --- onto \textit{kitobi}, and
that phrase is the subject of \textit{keldi}, ``arrived.''

\pause

The take-away: \say{strip the rightmost suffix and the syntax
disappears}. A parser that ignores it misses the parse. That is precisely
what motivates our \textbf{last-subword fusion} later in the talk. The
tree is a real example from our treebank.

% ------------------------------------------------------------
%  SLIDE 6 --- The gap and our approach
% ------------------------------------------------------------
\section*{Slide 6 --- The gap and our approach}
\timing{2:05}{2:35}

The core problem: existing Uzbek treebanks are small, and neural parsers
need \say{both} more data \say{and} better representations.

\begin{itemize}
   \item On the \textbf{data} side --- a new treebank, \textbf{UzUDT}, 684 gold sentences, the largest for Uzbek.
   \item On the \textbf{modelling} side --- two lightweight strategies: \textbf{last-subword fusion} and \textbf{cross-treebank training}.
\end{itemize}

\noindent
Crucially: \say{no heavy architecture changes}. This is the efficient,
low-resource-friendly path.

% ------------------------------------------------------------
%  SLIDE 7 --- Contributions
% ------------------------------------------------------------
\section*{Slide 7 --- Contributions}
\timing{2:35}{3:10}

Four contributions:

\begin{itemize}
   \item \textbf{UzUDT} --- built in INCEpTION, six annotators, full adjudication.
   \item \textbf{Last-subword fusion} --- a linguistically-motivated WordPiece-to-UD-token alignment.
   \item \textbf{Cross-treebank training} --- merging two genre-complementary Uzbek treebanks.
   \item A clean \textbf{factorial study} --- three configs $\times$ two data settings $=$ six runs --- plus released code and twelve checkpoints.
\end{itemize}

\noindent
Headline: best system \say{85.08 UPOS, 72.39 UAS, 63.81 LAS} --- and
cross-treebank alone adds \textbf{nine to eleven LAS points}.

% ------------------------------------------------------------
%  SLIDE 8 --- Related work
% ------------------------------------------------------------
\section*{Slide 8 --- Where this work fits}
\timing{3:10}{3:45}

Briefly --- the framework is \textbf{Universal Dependencies}; we build on
morphology-aware parsing for Turkic languages like Turkish. For Uzbek
specifically there was one small treebank, \say{Uzbek-UT}; we add
\say{UzUDT}.

\pause

Architecturally we use a standard neural pipeline --- Stanza-style, with
a \textbf{Dozat--Manning biaffine parser} --- and Uzbek encoders like
\textbf{TahrirchiBERT}. Our niche is \say{lightweight} robustness:
suffix-aware fusion plus cross-treebank data, not a new architecture.

% ------------------------------------------------------------
%  SLIDE 9 --- UzUDT
% ------------------------------------------------------------
\section*{Slide 9 --- UzUDT: a new gold Uzbek UD treebank}
\timing{3:45}{4:30}

UzUDT: \textbf{684 sentences}, about \textbf{7,800 tokens}, across fiction
and academic text. Annotated in INCEpTION by six people --- four
linguists and two NLP engineers --- with double annotation and
\say{full adjudication}. Full UD~v2 layers: UPOS, XPOS, lemma, features,
and dependency relations.

\pause

Inter-annotator agreement is strong --- \textbf{kappa of 95 on lemmas, 94
on UPOS, 91 on features}. So the labels are reliable. To our knowledge
this is the \say{largest publicly available Uzbek UD treebank}.

% ------------------------------------------------------------
%  SLIDE 10 --- Second treebank & data settings
% ------------------------------------------------------------
\section*{Slide 10 --- Second treebank \& data settings}
\timing{4:30}{5:05}

The second treebank is \textbf{UD\_Uzbek-UT} --- 500 sentences, news and
fiction, semi-automatic with manual correction.

\begin{itemize}
   \item Setting \textbf{.1} is UzUDT only --- 451 training sentences.
   \item Setting \textbf{.2} merges both --- 781.
\end{itemize}

\noindent
So merging nearly \say{doubles} the training data. That is a direct
test: \say{is raw data quantity the bottleneck?}

% ------------------------------------------------------------
%  SLIDE 11 --- Complementarity
% ------------------------------------------------------------
\section*{Slide 11 --- Why the treebanks are complementary}
\timing{5:05}{5:45}

These two treebanks cover \say{different syntax} --- which is exactly why
merging helps.

\begin{itemize}
   \item UT is news-heavy, so many \textbf{proper nouns} --- 5.2\% versus our 0.3\%.
   \item UzUDT has more \textbf{pronouns} and \textbf{four times} more adverbial clauses.
   \item And UT has \textbf{ten times} more light-verb compounds.
\end{itemize}

\noindent
So they fill each other's gaps --- merging is not just more data, it is
\say{broader} data.

% ------------------------------------------------------------
%  SLIDE 12 --- Pipeline overview
% ------------------------------------------------------------
\section*{Slide 12 --- Pipeline: from text to a labeled tree}
\timing{5:45}{6:30}

Here is the full pipeline, left to right. A UD-tokenised sentence goes
into a \say{frozen encoder} --- FastText or TahrirchiBERT. Then our
\textbf{fusion} step maps subwords to UD tokens. Then a \textbf{joint
tagger} predicts UPOS, XPOS and features. Then a \textbf{biaffine parser}
produces the labeled tree.

\pause

Two design choices matter. The encoder is \say{frozen} --- efficient.
And training is \say{sequential}: tagger first, then the parser runs on
\say{re-tagged} data --- the dashed arrow --- so there is no gold-tag
leakage at inference.

% ------------------------------------------------------------
%  SLIDE 13 --- Token representation
% ------------------------------------------------------------
\section*{Slide 13 --- Token representation: four channels}
\timing{6:30}{7:05}

Each token is a concatenation of four vectors: a static \textbf{word}
embedding, a \textbf{character}-LSTM vector, the \textbf{predicted UPOS}
embedding, and a \textbf{contextual} vector from BERT.

\pause

Two of the six runs differ here. In the \say{FastText} config, the word
channel is static 300-d and there is \say{no} contextual channel. In the
\say{TahrirchiBERT} config, that contextual channel is a \textbf{768-d
fused} vector --- and that fusion is our key contribution.

% ------------------------------------------------------------
%  SLIDE 14 --- Subword to token fusion (KEY)
% ------------------------------------------------------------
\section*{Slide 14 --- Subword$\to$token fusion: our key heuristic}
\timing{7:05}{8:05}

This is the heart of the method. BERT emits \say{WordPieces}, but UD
annotates \say{tokens} --- they do not line up. We must collapse several
subword vectors into one per token.

\pause

Our heuristic --- \say{last-subword}: take the vector of the \textbf{final}
WordPiece as the token's representation. For \textit{bolalarning} that is
\textit{\#\#ning}, the orange piece.

\pause

Why the last one? Because Uzbek is \say{suffixing} --- the rightmost
piece carries the \say{outermost} morphology. \textit{\#\#ning} literally
is the genitive case --- the bit the parser most needs. The baseline,
\textbf{mean pooling}, averages all pieces and \say{dilutes} exactly that
suffix signal. This one-line change is worth almost \textbf{four LAS
points}.

% ------------------------------------------------------------
%  SLIDE 15 --- Joint tagger
% ------------------------------------------------------------
\section*{Slide 15 --- Joint morphosyntactic tagger}
\timing{8:05}{8:45}

The tagger is a shared \textbf{BiLSTM} over the token vectors, feeding
\textbf{three softmax heads} --- UPOS, XPOS and morphological features ---
jointly. Joint tagging lets the three layers share evidence, which helps
when data is scarce.

\pause

And only the \say{predicted} UPOS embedding --- not the gold tag --- is
exported downstream to the parser. Again, \say{no leakage}.

% ------------------------------------------------------------
%  SLIDE 16 --- Biaffine parser
% ------------------------------------------------------------
\section*{Slide 16 --- Deep biaffine dependency parser}
\timing{8:45}{9:25}

Standard \textbf{Dozat--Manning} biaffine parser. The BiLSTM states ---
now including the predicted UPOS embedding --- are projected by MLPs into
\textbf{head} and \textbf{dependent} subspaces. A biaffine attention then
scores \say{every} head--dependent pair, and we decode a labeled tree
over the UD tokens.

\pause

The example on the right --- \textit{asarlar bizga eshigini ochib
beradi} --- shows subject, oblique, the \texttt{xcomp}, the object, and the
root.

% ------------------------------------------------------------
%  SLIDE 17 --- Experiment matrix
% ------------------------------------------------------------
\section*{Slide 17 --- Experiment matrix}
\timing{9:25}{10:00}

Six runs, fully factorial. Rows vary three things:

\begin{itemize}
   \item \textbf{E1 vs E2} --- the embedding: static FastText versus contextual BERT.
   \item \textbf{E2 vs E3} --- the fusion: our last-subword versus mean pooling.
   \item \textbf{.1 vs .2} --- the data: single versus cross-treebank.
\end{itemize}

\noindent
The green row, \textbf{E2.2}, is the full system: BERT, last-subword,
merged data.

% ------------------------------------------------------------
%  SLIDE 18 --- Main results
% ------------------------------------------------------------
\section*{Slide 18 --- Main results (test set)}
\timing{10:00}{11:00}

Tagging on the left, parsing on the right; the green row is our full
system, \textbf{E2.2}.

\begin{itemize}
   \item E2.2 wins \say{four of five} metrics --- best UPOS at \textbf{85.08}, best features, best UAS \textbf{72.39}, best LAS \textbf{63.81}.
   \item The only exception is \textbf{XPOS}, where mean pooling edges ahead --- a minor, fine-grained-tag effect.
\end{itemize}

\noindent
Notice the LAS column climbs steadily as we add \say{contextual
embeddings, then fusion, then data}.

% ------------------------------------------------------------
%  SLIDE 19 --- Three findings
% ------------------------------------------------------------
\section*{Slide 19 --- Three findings}
\timing{11:00}{12:00}

\begin{itemize}
   \item \textbf{1. Contextual beats static.} BERT over FastText gives nearly \textbf{five UPOS points} on merged data.
   \item \textbf{2. Fusion matters for parsing.} Last-subword over mean is \textbf{+3.76 LAS}. And the ordering is telling: mean-pooled BERT (60.05) actually \say{loses} to FastText (62.40); last-subword BERT (63.81) wins. A bad fusion wastes your encoder.
   \item \textbf{3. Cross-treebank is the biggest lever.} \textbf{+9.62 LAS} for BERT, \textbf{+11.16} for FastText --- just from merging.
\end{itemize}

\noindent
For low-resource parsing, \say{broader data beats clever modelling}.

% ------------------------------------------------------------
%  SLIDE 20 --- Error analysis
% ------------------------------------------------------------
\section*{Slide 20 --- Error analysis (100 test sentences, E2.2)}
\timing{12:00}{12:45}

We hand-checked 100 test sentences. The errors concentrate in
\say{attachment}:

\begin{itemize}
   \item \textbf{Prepositional / oblique attachment} --- 28\%.
   \item \textbf{Coordination scope} --- 22\%, then subordinate clauses and compounds.
\end{itemize}

\noindent
The dominant confusion is \textbf{\texttt{obl} versus \texttt{nmod}}, and
coordination scope. The tree on the right shows a missed \say{light-verb
construction} --- gold \texttt{compound:lvc} in blue, the model's wrong
arcs dashed in red. Rare constructions are exactly what the smaller
treebank under-covers.

% ------------------------------------------------------------
%  SLIDE 21 --- Conclusion
% ------------------------------------------------------------
\section*{Slide 21 --- Conclusion}
\timing{12:45}{13:30}

Three takeaways:

\begin{itemize}
   \item We released \textbf{UzUDT}, the largest gold Uzbek UD treebank.
   \item \textbf{Cross-treebank training} is the dominant factor --- up to eleven LAS points.
   \item \textbf{Last-subword fusion} is a cheap, linguistically-motivated win --- almost four LAS points.
\end{itemize}

\noindent
Best system: \say{85.08 / 72.39 / 63.81}. Limitations: small dev sets
and some genre mismatch. Next: \textbf{BERTbek} and
\textbf{BERT-plus-FastText fusion}. Everything --- data and twelve
checkpoints --- is on Hugging Face, code on GitHub.

% ------------------------------------------------------------
%  SLIDE 22 --- Thank you
% ------------------------------------------------------------
\section*{Slide 22 --- Thank you}
\timing{13:30}{14:00}

That brings me to the end. I would like to thank the open-source teams
behind \textbf{Stanza, TahrirchiBERT, and Universal Dependencies}, and
the UzUDT annotation team.

\pause

Thank you very much --- I am happy to take your questions.
(Contact: \texttt{s.matlatipov@nuu.uz}.)

% ------------------------------------------------------------
%  Appendix backups (only if asked)
% ------------------------------------------------------------
\newpage
\section*{Appendix backups --- only if asked}

\begin{itemize}
   \item \textbf{A1 --- Parser training dynamics.} Dev LAS/UAS over training; the merged setting converges higher and more stably than UzUDT-only.
   \item \textbf{A2 --- Tagger training dynamics.} UPOS / XPOS / UFeats / AllTags F1 across the six runs; E2.2 leads on most curves.
   \item \textbf{A3 --- A correctly parsed sentence.} A clean example the full system gets right end-to-end.
\end{itemize}

% ------------------------------------------------------------
%  Q & A prep --- printed as a back page for the speaker only
% ------------------------------------------------------------
\newpage
\section*{Likely Q\&A --- quick cues}

\noindent\textbf{Q. Why freeze the encoder instead of fine-tuning?}\\
\textbf{Efficiency and stability}. On fewer than a thousand sentences,
fine-tuning a Transformer overfits; freezing gives a strong, reproducible
baseline at a fraction of the cost.

\vspace{0.4em}
\noindent\textbf{Q. Why last-subword pooling rather than mean or attention?}\\
It is \textbf{motivated by Uzbek morphology} --- the rightmost subword
carries the outermost, syntactically decisive suffix. Empirically it is
\textbf{+3.76 LAS} over mean pooling, with no extra parameters.

\vspace{0.4em}
\noindent\textbf{Q. Is last-subword fusion Uzbek-specific?}\\
It targets \textbf{suffixing, agglutinative} morphology, so it should
transfer to other Turkic languages --- Turkish, Kazakh, Uyghur. Testing
that is future work.

\vspace{0.4em}
\noindent\textbf{Q. Why not mBERT, XLM-R, or UDify?}\\
We deliberately compare \textbf{lightweight, Uzbek-specific} encoders for
a fair low-resource baseline. Multilingual and cross-lingual transfer is
a natural --- but separate --- next step.

\vspace{0.4em}
\noindent\textbf{Q. The dev sets are tiny (45 / 78) --- isn't that noisy?}\\
\textbf{Yes, and we state it as a limitation.} But the headline data
effect --- \textbf{nine to eleven LAS points} --- is far larger than any
plausible dev-set noise, so the conclusion is safe.

\vspace{0.4em}
\noindent\textbf{Q. Why does XPOS regress under last-subword?}\\
Mean pooling slightly \textbf{smooths fine-grained tags}. The net effect
across the other four metrics still favours last-subword, and it is our
best overall system (four of five metrics).

\vspace{0.4em}
\noindent\textbf{Q. Does merging help only because of size, or genre too?}\\
\textbf{Both.} Merging doubles the data \emph{and} the treebanks are
genre-complementary --- UT brings proper nouns and light-verb compounds,
UzUDT brings pronouns and adverbial clauses.

\vspace{0.4em}
\noindent\textbf{Q. Will the code and checkpoints be released?}\\
\textbf{Yes} --- the UzUDT data and \textbf{twelve checkpoints} are on
Hugging Face, and the full code is on GitHub.

\end{document}
